\hiddensection{Kurzfassung}
Die vorliegende Arbeit stellt einen neuen Ansatz für die Erzeugung glaubwürdiger \acp{npc} in Videospielen vor. Der Fokus liegt auf dem Entscheidungsprozess von intelligenten Agenten für die Darstellung eines authentischen Verhaltens. Anstelle von klassischen Methoden wie \Fachbegriff{State Machines} basiert der entwickelte Ansatz auf \Fachbegriff{Reinforcement Learning} und \Fachbegriff{Utility Based Agents}. Dabei handelt es sich um  bedürfnisorientierte Agenten, die über menschliche Attribute verfügen. Sie lernen ihre Entscheidungen so zu wählen, dass ihre Bedürfnisfunktion maximiert wird.

Als Voraussetzung für das Verfahren werden Grundlagen zu intelligenten Agenten sowie im Machine Learning-Bereich mit starkem Fokus auf Reinforcement Learning geschaffen. In der Spieleindustrie etablierte Technologien für die Entwicklung von \acp{npc} werden diskutiert, um einen direkten Vergleich mit dem eigenen Verfahren herstellen zu können.
Der Hauptteil der Arbeit besteht aus einer Implementierung des Verfahrens anhand eines bewohnten Dorfsystems als Umgebung für lernende Agenten mithilfe der Game Engine Unity. Die entstandene Architektur sowie ihre einzelnen Komponenten werden ausführlich diskutiert.

Ferner werden darauf aufbauend zwei zusätzliche Ansätze entwickelt, um verschiedene Methoden gegeneinander abzuwägen und Vergleiche zu ziehen. Genauer werden eine regelbasierte Vorgehensweise sowie \acp{goap} verwendet.

Schlussendlich folgt ein Überblick über mögliche weiterführende Entwicklungen und ein Fazit zu den Implementierungen mit einem Fokus auf die Beurteilung der Glaubwürdigkeit der \acp{npc} und die Schwierigkeiten der jeweiligen Technologien. Dabei werden der Erfolg des entwickelten Ansatzes auch im Hinblick auf sowohl vergleichbare als auch unterschiedliche relevante Verfahren bewertet und seine Schwierigkeiten sowie Vorteile herausgestellt. Zuletzt liefert ein Ausblick Informationen über mögliche zukünftige Forschung.




\newpage




\hiddensection{Abstract}
This study introduces a new approach towards the creation of believable \FachbegriffMitAbk{non-playable-characters}{npcs} in videogames. Its focus is set on improving the internal decision process utilized in intelligent agents to induce a more authentic environment for the player. A Machine Learning approach is applied using \Fachbegriff{Reinforcement Learning} rather than classical methods like \Fachbegriff{Finite State Machines}. Furthermore the agent is encouraged to display humanlike behavior by maximizing need-based utility functions.

As a preliminary requirement for the proposed method, fundamental knowledge is provided by illustrating underlying concepts such as intelligent agents and Reinforcement Learning. Well-established technologies for the definition of NPC behavior is discussed in order to create a direct comparison to the designed process. The main part of the thesis comprises the implementation of said method by developing a village simulation as an environment for learning agents using the game engine \Fachbegriff{Unity}. A general outline of the architecture as well as a detailed description of the various components employed in this system is then conducted.

Moreover, two additional methods are applied to create a direct comparison to draw conclusions and weigh the benefits and tracks of the different approaches. Specifically this includes a rule-based approach as well as \Fachbegriff{Goal Oriented Action Planning} (GOAP).
Finally an overview about further improvements and a conclusion about the generated approaches is given. The latter includes an assessment of the believability of the respective methods and their difficulies as well as advantages of the underlying technology determining the success of the method. An outlook then provides an overview about possible further research.


% alter text: 

%Authentic behavior of \FachbegriffMitAbk{non-playable-characters}{npcs} is paramount for the creation of immersive and plausible virtual worlds.
%The believability and atmosphere of a player's environment is largely influenced by the actions of the NPCs inhabiting it and their achieved degree of humanlike behavior. This includes not only interactable NPCs but also background characters that are often used to populate a world by enhancing its authenticity and making it feel more alive. Those mentioned features were continuously improved since the popularization of video games - however when observing specific NPCs separately, it becomes apparent that they often do not convey humanlike behavior but instead seem predictable and hardly autonomous.

%This thesis focuses on finding a possible solution for the aforementioned problem while focusing on meaningful actions that NPCs can take and connected decision processes rather than visual aspects. An approach for a \Fachbegriff{Reinforcement Learning} based implementation with \textit{utility-based} or \textit{need-based} agents is suggested. A Unity prototype is developed to demonstrate the agents in action. Their purpose is to learn authentic behavior themselves by utilizing their human-based need values and therefore creating humanlike manners and a valid dynamic daily routine.
%As a prerequisite, fundamental Machine Learning principles, specifically in the field of \textit{Reinforcement Learning} are examined. The focus lies on implementation possibilities in the game engine \Fachbegriff{Unity}.

%Further more an introduction into the history and current state of \Fachbegriff{intelligent agents} or \Fachbegriff{game AI} in videogames is conducted.
%Based on this knowledge additional approaches utilizing different popular gaming AI technologies are discussed and developed to compare the results and draw conclusions.

%Finally an outlook is given into possible further research. The results are evaluated and conclusions are drawn to estimate the authenticity of the generated NPCs as well as the success of the developed approach by comparing it to other standard technologies.

\newpage